% Generated draft: requires mathematical and source review.
\documentclass[11pt]{article}
\usepackage[margin=25mm]{geometry}
\usepackage{fontspec}
\setmainfont{Latin Modern Roman}
\usepackage{amsmath,amssymb,mathtools,unicode-math}
\setmathfont{Latin Modern Math}
\usepackage{enumitem,needspace,xurl,hyperref}
\setlength{\parindent}{0pt}
\setlength{\parskip}{5pt}
\setlength{\emergencystretch}{3em}
\newcommand{\N}{\mathbb{N}}
\newcommand{\Z}{\mathbb{Z}}
\newcommand{\Q}{\mathbb{Q}}
\newcommand{\R}{\mathbb{R}}
\newcommand{\C}{\mathbb{C}}
\newcommand{\F}{\mathbb{F}}
\newcommand{\A}{\mathbb{A}}
\newcommand{\PP}{\mathbb{P}}
\newcommand{\E}{\mathbb{E}}
\newcommand{\eps}{\varepsilon}
\newcommand{\dd}{\,\mathrm{d}}
\begin{document}
\section*{UnsolvedMath ID: 30005705}
\textbf{Source problem code:} catalogue entry
\subsection{Definitions and mathematical statement}
A totally disconnected locally compact (TDLC) group is a Hausdorff locally compact topological group whose connected components are singletons. It is compactly generated if some compact subset generates it as an abstract group.

For a TDLC group \(G\), let \(\mathcal D_{\Q}(G)\) be the category of \(\Q\)-vector spaces with a linear \(G\)-action such that every vector has an open stabilizer; equivalently, these are continuous representations on discrete \(\Q\)-vector spaces. Interpret the rational discrete cohomological dimension as
\[
\operatorname{cd}_{\Q}G
=\operatorname{pd}_{\mathcal D_{\Q}(G)}\Q,
\]
where \(\Q\) has trivial \(G\)-action and projective dimension is infinite if no finite projective resolution exists. This interpretation follows the supplied literature assessment's specification of rational discrete cohomological dimension.

For compactly generated groups \(G,H\), choose compact symmetric generating sets \(S,T\) containing the identity and equip the groups with their word metrics \(d_S,d_T\). They are quasi-isometric if there exist a map \(f:G\to H\) and constants \(A\geq 1\), \(B,R\geq 0\) such that
\[
A^{-1}d_S(g,g')-B
\leq d_T(f(g),f(g'))
\leq A d_S(g,g')+B
\quad(g,g'\in G),
\]
and every point of \(H\) has distance at most \(R\) from \(f(G)\). This condition is independent of the chosen compact generating sets; the map is not required to be continuous or a homomorphism.

The first question in the supplied clean statement is: for every pair of compactly generated TDLC groups \(G,H\), does
\[
G\simeq_{\mathrm{qi}}H,
\qquad \operatorname{cd}_{\Q}G<\infty,
\qquad \operatorname{cd}_{\Q}H<\infty
\quad\Longrightarrow\quad
\operatorname{cd}_{\Q}G=\operatorname{cd}_{\Q}H
\]
hold?

For a simple graph \(\Lambda=(V,E)\), the associated right-angled Artin group is
\[
A_{\Lambda}
=\left\langle v\ (v\in V)\ \middle|\
vw=wv\text{ whenever }\{v,w\}\in E\right\rangle.
\]
The second question asks for a suitable or useful TDLC analogue of these groups. No required construction, axioms, or criterion of usefulness is supplied, so this part remains a research request rather than a uniquely specified mathematical proposition. The supplied dated assessment reports a concrete topological right-angled Artin group construction in a 2026 preprint, while reporting that the first question remains unsettled; this is supplied curation, not a conclusion established here.

The raw original excerpt also contains adjacent material whose inclusion in this permanent record's intended target is unclear. Under the name Rudradip Biswas, it quotes a result labelled \([1,\ \text{Theorem }19.1]\): if \(\Gamma\in\mathrm{LHF}\), described there as locally in Kropholler's hierarchy with all finite groups as the base class, and \(\Gamma\) is of type \(\mathrm{FP}_{\infty}(\Z)\), then \(\Gamma\) has only finitely many conjugacy classes of finite elementary abelian \(p\)-subgroups. Here an elementary abelian \(p\)-group is \((\Z/p\Z)^r\) for a prime \(p\) and a nonnegative integer \(r\). The excerpt does not explicitly quantify \(p\), and its reference \([1]\) is not identified. Its Question 2 asks whether a TDLC version of that result can be formulated, with a remark suggesting consideration of type \(\mathrm{FP}_{\infty}(\Q)\) or \(\mathrm F_{\infty}\). For a discrete group, type \(\mathrm{FP}_{\infty}(R)\) means that the trivial \(R\Gamma\)-module \(R\) has a projective resolution finitely generated in every degree; type \(\mathrm F_{\infty}\) means that a classifying space has finite skeleton in every dimension. The excerpt does not specify the intended TDLC replacements for these conditions or the conclusion.

Under Ian Leary and Ilaria Castellano, the same raw excerpt gives Question 5: does every \(\sigma\)-compact TDLC group embed into a compactly generated TDLC group? Here \(\sigma\)-compact means a countable union of compact subsets. The excerpt does not specify whether the embedding must be a topological embedding or satisfy additional image conditions. These adjacent questions are preserved as contextual source material, without identifying them with the two questions of the clean statement.
\subsection{Short English statement}
For quasi-isometric compactly generated TDLC groups with finite rational discrete cohomological dimension, must those dimensions agree? The record also asks for a useful TDLC analogue of right-angled Artin groups, without specifying what properties that analogue must have. Its supplied dated assessment labels the compound record open and reports a construction relevant to the second question. Additional questions in the raw excerpt have uncertain scope.
\subsection{Sources}
\begin{itemize}
\item Local UnsolvedMath registry: permanent ID \(30005705\), title \emph{Quasi-Isometry Invariance and RAAG Analogues for TDLC Groups}, exact source problem code \texttt{catalogue entry}, in the Oberwolfach Reports Open Problems collection. The supplied record labels the status \texttt{open} and the statement status \texttt{corrected\_verified}; its assertion of checking against the official report is a supplied verification claim, not an independent verification here. No proposer is specified for the principal questions.
\item Supplied source citation: \emph{Mini-Workshop: Homological aspects of TDLC-groups (2024)}, Oberwolfach Reports, pp. 3103--3138, DOI \texttt{10.4171/owr/2023/55}. \url{https://doi.org/10.4171/owr/2023/55}. The record separately gives proposed year 2023 and the tag \texttt{owr-year:2023}; the dating discrepancy is unresolved.
\item Supplied literature source, labelled ``Oberwolfach report'': \url{https://ems.press/content/serial-article-files/48166}.
\item Supplied literature source, attributed to Castellano--Marchionna--Nucinkis--Santos Rego: \url{https://arxiv.org/abs/2605.28963}. The supplied assessment describes this as a 2026 preprint providing a concrete topological right-angled Artin group construction, without settling quasi-isometry invariance of rational discrete cohomological dimension.
\item Literature assessment checked at \texttt{2026-08-22}: the supplied status is open. This is supplied dated curation derived from the OWR record, not independently verified; its time-sensitive status and cited claims require independent checking.
\item The raw excerpt attributes adjacent material to Rudradip Biswas and to Ian Leary and Ilaria Castellano. Its citation \([1,\ \text{Theorem }19.1]\) has no accompanying bibliographic identification in the supplied record.
\end{itemize}
\end{document}
